\section{Erd\H{o}s Problem 269 --- Round 2}

\subsection*{1) ROUND-2 OBJECTIVE}
Path \textbf{(C)}: the finite-set problem remains open, so the most promising strict advance is to isolate the multiplicity obstruction identified in Round~1, turn it into a precise irrationality criterion, and then prove a rigorous irrationality theorem for a broad infinite-prime-set regime. In addition, I explicitly re-check the unverified forum proof strategy and repair only the part that survives adversarial scrutiny.

\subsection*{2) Round-1 FOUNDATION USED}
I use the following Round~1 items without re-proving them.

\begin{itemize}
\item The sequence of distinct denominators can be studied by grouping equal values of
\[
L_n=[a_1,\dots,a_n].
\]
Round~1 already identified the main obstruction: the original series is a weighted version
\[
\sum_i \frac{c_i}{d_i}
\]
with multiplicities $c_i$, and these multiplicities are what prevent the simple Cantor-series tail argument from resolving the finite case.
\item Round~1 also proved convergence of the series, so all rearrangements into grouped denominator levels are legitimate.
\end{itemize}

\subsection*{3) NEW INSIGHT / TOOL (ROUND-2)}
The new tool is a prime-power frontier decomposition.

Let $P$ be any set of primes, and let
\[
1<q_1<q_2<\cdots
\]
be the increasing sequence of all prime powers $p^r$ with $p\in P$ and $r\ge 1$. Define
\[
q_0:=1,
\qquad
x_i:=\mathrm{lcm}(q_1,\dots,q_i)\quad (i\ge 1),
\qquad x_0:=1.
\]
For each $i\ge 0$, let $m_i$ be the number of $P$-smooth integers in the half-open interval $[q_i,q_{i+1})$.

Then the original series becomes
\[
\sum_{n\ge 1}\frac1{L_n}
=
\sum_{i\ge 0}\frac{m_i}{x_i}.
\]
The key Round-2 criterion is:

\medskip
\noindent
\textbf{Criterion.}
If $P$ is infinite and
\[
\frac{m_i}{q_i}\to 0,
\]
then the grouped series $\sum_{i\ge 0} m_i/x_i$ is irrational.
\medskip

This is strictly stronger than anything obtained in Round~1, and it pinpoints a single sharp obstruction: one only needs sufficiently small multiplicity relative to the prime-power frontier.

\subsection*{4) ATTACK PLAN (ROUND-2)}
The exact remaining gap after Round~1 was: repeated denominators can be too numerous for the simple tail $<1/d_i$ argument.

The new plan is:

\begin{itemize}
\item encode denominator changes by the ordered prime powers $q_i$;
\item show that the tail after a \emph{new prime} enters the denominator is an integer multiple of $1/x_N$ under a rationality assumption;
\item show that this integer tends to $0$ whenever $m_i=o(q_i)$;
\item deduce irrationality in every situation where one can actually prove $m_i=o(q_i)$.
\end{itemize}

This directly overcomes the Round~1 obstacle in all regimes where the multiplicities are genuinely sparse relative to the frontier.

\subsection*{5) WORK (ROUND-2)}

\paragraph{Lemma 5.1.}
If $q_{i+1}=p^r$, then
\[
x_{i+1}=p\,x_i.
\]
Equivalently, if $\alpha(q_{i+1})$ denotes the prime base of $q_{i+1}$, then
\[
x_{i+1}=\alpha(q_{i+1})x_i.
\]

\emph{Proof.}
Since the $q_j$ are the distinct prime powers coming from $P$, if $q_{i+1}=p^r$ then $p^{r-1}$ already occurred earlier but $p^r$ did not. Thus the $p$-adic valuation of the lcm increases by exactly $1$, while all other prime-adic valuations stay unchanged. Hence $x_{i+1}=p x_i$. 

\paragraph{Lemma 5.2.}
For every $i\ge 0$,
\[
\sum_{n\ge 1}\frac1{L_n}=\sum_{i\ge 0}\frac{m_i}{x_i}.
\]

\emph{Proof.}
Fix $i\ge 0$. For every $P$-smooth integer $u$ with $q_i\le u<q_{i+1}$, all prime powers from $P$ not exceeding $u$ are precisely $q_1,\dots,q_i$, so
\[
\mathrm{lcm}\{v\in\mathbf Z_{>0}: v\le u,\ v\text{ is }P\text{-smooth}\}=x_i.
\]
Exactly $m_i$ values of $u$ lie in $[q_i,q_{i+1})$, so grouping equal denominators gives the formula. 

\paragraph{Lemma 5.3.}
Let $p_1$ be the smallest prime in $P$. Then for every $i\ge 0$,
\[
q_{i+1}\le p_1 q_i.
\]
Consequently,
\[
q_i\le p_1^{\,i-N-1}q_{N+1}\qquad (i\ge N+1).
\]

\emph{Proof.}
Given $q_i$, choose the unique integer $s$ with
\[
p_1^{s-1}\le q_i < p_1^s.
\]
Then $p_1^s$ is a prime power from $P$, it is $>q_i$, and $p_1^s\le p_1 q_i$. Since $q_{i+1}$ is the next prime power after $q_i$, we get $q_{i+1}\le p_1 q_i$. Iterating gives the second claim. 

\paragraph{Proposition 5.4 (irrationality criterion from sparse multiplicities).}
Assume that $P$ is infinite and that
\[
\frac{m_i}{q_i}\to 0.
\]
Then
\[
\sum_{n\ge 1}\frac1{L_n}
\]
is irrational.

\emph{Proof.}
Write
\[
S:=\sum_{n\ge 1}\frac1{L_n}=\sum_{i\ge 0}\frac{m_i}{x_i}.
\]
Assume for contradiction that $S=a/b\in\mathbf Q$ in lowest terms.

Because $P$ is infinite, infinitely many of the $q_j$ are actual primes. Choose $N$ so that $q_{N+1}$ is a prime. Define the tail
\[
T_N:=\sum_{i\ge N+1}\frac{m_i}{x_i}.
\]
Since $x_i\mid x_N$ for $i\le N$, we have
\[
I_N:=b x_N T_N
=
ax_N-bx_N\sum_{i=0}^{N}\frac{m_i}{x_i}
\in \mathbf Z_{>0}.
\]
So $I_N$ is a positive integer.

Now let
\[
M_N:=\sup_{i\ge N+1}\frac{m_i}{q_i}.
\]
Since $m_i/q_i\to 0$, we have $M_N\to 0$ as $N\to\infty$.

Using Lemma~5.1,
\[
\frac{x_N}{x_i}
=
\frac1{\alpha(q_{N+1})\alpha(q_{N+2})\cdots \alpha(q_i)}.
\]
Hence
\[
I_N
=
 b\sum_{i\ge N+1}
 \frac{m_i}{\alpha(q_{N+1})\alpha(q_{N+2})\cdots \alpha(q_i)}
\le
 b M_N
 \sum_{i\ge N+1}
 \frac{q_i}{\alpha(q_{N+1})\alpha(q_{N+2})\cdots \alpha(q_i)}.
\]
Because $q_{N+1}$ is a prime, $\alpha(q_{N+1})=q_{N+1}$. By Lemma~5.3,
\[
\frac{q_i}{\alpha(q_{N+1})\cdots \alpha(q_i)}
\le
\frac{p_1^{\,i-N-1}}{\alpha(q_{N+2})\cdots \alpha(q_i)}.
\]
Let $p_2$ be the second-smallest prime in $P$, and let
\[
n_i:=\#\{j\in\{N+2,\dots,i\}: \alpha(q_j)>p_1\}.
\]
Since each factor $\alpha(q_j)$ is either $p_1$ or at least $p_2$,
\[
\alpha(q_{N+2})\cdots \alpha(q_i)
\ge
p_1^{\,i-N-1-n_i}p_2^{n_i},
\]
so
\[
\frac{q_i}{\alpha(q_{N+1})\cdots \alpha(q_i)}
\le
\left(\frac{p_1}{p_2}\right)^{n_i}.
\]
Therefore
\[
I_N\le b M_N \sum_{i\ge N+1}\left(\frac{p_1}{p_2}\right)^{n_i}.
\]
It remains to bound the last series uniformly in $N$.

Set
\[
\ell:=\left\lceil \frac{\log p_2}{\log p_1}\right\rceil.
\]
For every $x>0$, the next power of $p_2$ after $x$ is at most $p_2 x\le p_1^{\ell}x$. Thus every interval $(x,p_1^{\ell}x]$ contains a power of $p_2$. In particular, it is impossible to have $\ell+1$ consecutive terms of the prime-power sequence $(q_i)$ all with base prime $p_1$: the first and last of those terms would differ by a factor at most $p_1^{\ell}$, but there would have to be a power of $p_2$ strictly between them. Hence $n_i$ can stay constant for at most $\ell+1$ consecutive indices. Therefore, for each $k\ge 0$,
\[
\#\{i\ge N+1: n_i=k\}\le \ell+1.
\]
So
\[
\sum_{i\ge N+1}\left(\frac{p_1}{p_2}\right)^{n_i}
\le
(\ell+1)\sum_{k\ge 0}\left(\frac{p_1}{p_2}\right)^k
=:C(P)<\infty,
\]
with $C(P)$ independent of $N$.

Hence
\[
0<I_N\le b C(P) M_N\to 0.
\]
But $I_N$ is a positive integer, contradiction. Therefore $S$ is irrational. 

\paragraph{Corollary 5.5.}
Assume $P$ is infinite and one of the following holds.

\begin{itemize}
\item[(i)] $q_{i+1}-q_i=o(q_i)$ for the ordered prime-power frontier;
\item[(ii)] $\sum_{p\notin P}1/p=\infty$.
\end{itemize}
Then
\[
\sum_{n\ge 1}\frac1{L_n}
\]
is irrational.

\emph{Proof.}
In case (i),
\[
m_i\le q_{i+1}-q_i=o(q_i),
\]
so Proposition~5.4 applies.

In case (ii), let $h_i:=q_{i+1}-q_i$. Every integer counted by $m_i$ lies in an interval of length $h_i$ and has no prime factor outside $P$. Brun's upper-bound sieve gives
\[
m_i\ll h_i\prod_{\substack{p\le h_i\\ p\notin P}}\left(1-\frac1p\right).
\]
If $h_i=o(q_i)$ then we are done. Otherwise there is some $\varepsilon>0$ along an infinite subsequence with $h_i\ge \varepsilon q_i$, so $h_i\to\infty$ there; because $\sum_{p\notin P}1/p=\infty$, the Euler product above tends to $0$, and thus along that subsequence
\[
m_i=o(h_i)=o(q_i).
\]
Therefore $m_i=o(q_i)$ in all cases, and Proposition~5.4 applies. 

\paragraph{Round-2 consequence.}
The proposition and corollary give a complete irrationality proof for a broad infinite-$P$ regime, including for example the case where $P$ is the set of all primes.

\subsection*{6) ADVERSARIAL VERIFICATION}
There is one important correction to make relative to the unverified forum proof sketch.

\begin{itemize}
\item The step ``$\sum_{p\notin P}1/p<\infty$ implies every interval $(x,(1+\varepsilon)x]$ eventually contains a prime of $P$'' is not justified by the argument sketched on the forum page. Convergence of the complementary reciprocal sum alone does not obviously force this multiplicative-syndetic property.
\item The proof written above therefore \emph{does not} claim irrationality for every infinite $P$. What it proves is the precise conditional theorem in Proposition~5.4, together with the rigorous Corollary~5.5 covering two broad classes of infinite sets.
\item Every use of positivity and integrality in Proposition~5.4 has been checked: $I_N>0$ because the tail is positive, and $I_N\in\mathbf Z$ because $x_i\mid x_N$ for all $i\le N$.
\item The boundedness of the auxiliary series was also checked carefully: the only input is that powers of $p_2$ occur regularly on the multiplicative scale, and Lemma~5.3 converts this into a uniform bound on the number of indices with a fixed value of $n_i$.
\end{itemize}

So the new argument is rigorous as stated, but it still leaves open the genuinely difficult remaining infinite-$P$ cases, and of course the original finite-$P$ problem.

\subsection*{7) FINAL (EXACTLY ONE)}
\textbf{UNRESOLVED (BUT STRICTLY ADVANCED)}

I proved a new irrationality criterion for the grouped series and obtained a rigorous irrationality theorem for all infinite prime sets $P$ satisfying $m_i=o(q_i)$, in particular for every infinite $P$ with either $q_{i+1}-q_i=o(q_i)$ or $\sum_{p\notin P}1/p=\infty$. I also identified a concrete gap in the previously cited unverified ``all infinite $P$'' argument. The original finite-$P$ problem remains open.

\subsection*{8) COMPLETION ESTIMATE}
COMPLETION: 45\%

\subsection*{9) REFERENCES}
\begin{itemize}
\item T. F. Bloom, \emph{Erd\H{o}s Problem \#269}, checked 2026-03-07.
\item T. F. Bloom, \emph{Discussion thread for Erd\H{o}s Problem \#269}; in particular the comment of Steve Fan (11 Oct 2025), checked 2026-03-07. Used as a heuristic starting point only; the Round-2 proof above is the corrected version actually verified here.
\end{itemize}

